\section{来源审计：本讲真正要改变的是 RL 的问题表述}

本讲介绍的不是一个“把 Transformer 当作 policy network”的局部架构替换，而是一种更根本的重新表述：把固定数据集中的轨迹写成 token 序列，把期望回报写成条件，再用 causal Transformer 完成动作预测。这样做的价值不只在模型结构，还在于它把训练目标从不稳定的动态规划更新转向了熟悉的监督学习损失；相应地，它也继承了序列模型对数据支持、context length 和条件外推的限制。

\subsection{来源、时间边界与日期冲突}

这份讲义以 Stanford Online 官方视频、官方手动英文字幕、视频中恢复的 24 张教学 slide、Decision Transformer 原始论文与官方代码为主要证据。Stanford CS25 V1 课程页把本场列在 2021 年 10 月 11 日，论文属于 NeurIPS 2021；但视频中的标题页写着 ``October 11, 2022''，而公开视频上传日期又是 2022 年 7 月 13 日。三者互相冲突，因此本文把课程日程和论文出版记录作为历史锚点，同时明确保留画面中的日期异常，不为来源自动“修正事实”。

\lecturefigure{slide-01-title.jpg}{讲座标题页：Decision Transformer 把 reinforcement learning 表述为 sequence modeling。}{Stanford Online 官方视频 00:06；画面日期与课程页、上传日期存在冲突，见来源审计。}

\begin{warningbox}{来源冲突不能靠猜测消失}
当课程页、视频元数据和 slide 文字发生冲突时，正确做法是记录冲突、说明采用哪一个时间锚点，并把技术结论限制在不依赖该冲突的部分。这里可以确认的是：核心论文为 NeurIPS 2021 工作，公开视频由 Stanford Online 发布；无法仅凭标题页断言实际录制年份。
\end{warningbox}

\subsection{本讲的三个核心问题}

阅读全文时应持续追踪三个问题。第一，为什么传统 RL 的训练和扩展经常比 language modeling 更脆弱？第二，return-to-go 如何把“我想得到多高回报”变成一个可输入的控制条件？第三，实验中哪些提升来自序列建模本身，哪些仍依赖离线数据的覆盖、筛选和 benchmark 设计？这三个问题共同决定 Decision Transformer 是概念突破、工程工具，还是被过度概括的口号。

\begin{importantbox}{本讲学习目标}
完成本讲后，读者应能：
\begin{enumerate}
  \item 写出 MDP、return-to-go 与 Decision Transformer token sequence；
  \item 解释 causal mask、timestep embedding、context length 与 non-Markov conditioning；
  \item 从训练到 rollout 手算一次 target return 的更新；
  \item 区分 requested return、realized return、logged-data support 与 extrapolation；
  \item 解释 percent behavior cloning、CQL、sparse reward、credit assignment 和 critic 实验；
  \item 说明为什么 offline RL 的成功不能直接推出 online RL 已经解决。
\end{enumerate}
\end{importantbox}

\subsection{本章小结}

本讲的历史来源存在可见日期冲突，但技术主线清楚：以 NeurIPS 2021 的 Decision Transformer 为核心，讨论如何将 offline RL 重构为 conditional sequence modeling。后续每个结果都必须同时检查模型机制、数据支持和评价方式。

